\documentclass[11pt]{article}

\usepackage[margin=1.1in]{geometry}
\usepackage{amsmath}
\usepackage{amssymb}
\usepackage{amsthm}
\usepackage{mathtools}
\usepackage[expansion=false]{microtype}
\usepackage{hyperref}
\hypersetup{colorlinks=true, linkcolor=black, citecolor=black, urlcolor=blue}

\emergencystretch 2.5em

\newtheorem{theorem}{Theorem}
\newtheorem{proposition}[theorem]{Proposition}
\newtheorem{lemma}[theorem]{Lemma}
\newtheorem{corollary}[theorem]{Corollary}
\theoremstyle{definition}
\newtheorem{definition}[theorem]{Definition}
\newtheorem{convention}[theorem]{Convention}
\newtheorem{obligation}{Obligation}
\theoremstyle{remark}
\newtheorem{remark}[theorem]{Remark}
\newtheorem{example}[theorem]{Example}

% --- house macros -----------------------------------------------------------
% The calculus is never written with the Greek letter that follows pi.
\newcommand{\rhoc}{\textsf{rho}}
\newcommand{\Proc}{\mathsf{Proc}}
\newcommand{\Name}{\mathsf{Name}}
\newcommand{\Pair}{\mathsf{Pair}}
\newcommand{\pzero}{\mathsf{0}}
\newcommand{\quo}[1]{\ensuremath{@_{#1}}}
\newcommand{\drp}[1]{\ensuremath{*_{#1}}}
\newcommand{\dip}[2]{\ensuremath{\langle #1 , #2 \rangle}}
\newcommand{\negop}{\ensuremath{\ominus}}
\newcommand{\FN}[1]{\ensuremath{\mathsf{FN}(#1)}}
\newcommand{\sect}[1]{\ensuremath{\mathsf{sec}(#1)}}
\newcommand{\addr}[1]{\ensuremath{\mathsf{addr}(#1)}}
\newcommand{\phase}[1]{\ensuremath{\mathsf{ph}(#1)}}
\newcommand{\scong}{\ensuremath{\equiv}}
\newcommand{\nameq}{\ensuremath{\equiv_{N}}}
\newcommand{\red}{\ensuremath{\rightarrow}}
\newcommand{\Id}{\mathbb{1}}

\title{A Complex Namespace for the \rhoc\ Calculus\\[2mm]
\large Paired messages, two quotations, and an intertwining COMM}
\author{L.G. Meredith\\ \small F1R3FLY.io}
\date{Version 2 --- a record of the core calculus}

\begin{document}
\maketitle

\begin{abstract}
This note records, in one place, the core of a variant of the \rhoc\ calculus
designed so that the namespace carries the structure of a complex scalar field
rather than that of a free commutative monoid. Three changes are made to the
term language and one to the reduction relation. Messages carry \emph{pairs} of
processes, so that a message is a formal difference and negation is structural.
Names are quotations of pairs, so that negation is an operation on
\emph{addresses} and not merely a decoration on sends. There are two quotation
operators and, of necessity, two dropping operators; the two quotations are read
as a real and an imaginary sector. Communication \emph{intertwines} the sectors,
so that a communication event is multiplication by the imaginary unit, and
communication between participants of opposite polarity exchanges the components
of the received pair, so that opposite polarity contributes the sign. The note
states what has been settled, states the design forks that have not, and
enumerates the obligations the design incurs. It also records the evaluation
product against which the eventual pairing is defined, together with the pole
and the residuation that give the dagger, taking the pole to be the zero process
and the perp to be residuation to it. It proves almost nothing. Its
purpose is to fix notation and to make the open questions precise enough to be
worked on.
\end{abstract}

\tableofcontents

% ---------------------------------------------------------------------------
\section{What this note is}

\subsection{Scope}

The present document is a \emph{record}, not an argument. It writes down the
term language, the structural theory and the reduction relation of a variant
calculus that has emerged from a sequence of design decisions, together with
the reasons each piece is present. It does not develop the semantics that
motivated the design; that semantics --- an interpretation of the operations of
quantum mechanics over a reflective process calculus, with processes as
vectors, names as scalars, evaluated contexts as covectors and contexts as
operators --- is the subject of separate work, and is recalled here only far
enough to explain why each feature was added.

Two things are deliberately excluded. The note gives no metatheory: no subject
reduction, no congruence result, no decidability. And it gives no semantics: no
interpretation map, no inner product, no dagger beyond the one line needed to
say which operation the dagger is built from. Both are named as obligations in
Section~\ref{sec:obligations}.

\subsection{Status markers}

Throughout, claims carry one of three markers.

\begin{description}
\item[SETTLED] The design decision has been taken and the consequences worked
  through. It may still be wrong, but it is not open.
\item[FORK] Two or more readings are consistent with what has been settled, and
  the choice has not been made. The alternatives are stated.
\item[OPEN] An obligation the design incurs and does not discharge.
\end{description}

\subsection{Revision note}

Version 1 of this note had receipts bind a pair of name variables and stated the
communication rule in two clauses, one per relative polarity. Both features are
withdrawn. Negation on names lifts uniquely to processes
(Proposition~\ref{prop:lift}), the whole received pair is deposited into a single
bound variable, and the polarity clause becomes a factor in one rule. The change
is not cosmetic: requiring reduction to be equivariant for the lifted negation
\emph{selects} the deposit convention, and version 1's convention fails that
test (Remark~\ref{rem:selection}). The surviving pair sort occurs only on
messages and inside quotations.

\subsection{Notational rules}

The calculus is called the \rhoc\ calculus and is never written with the Greek
letter that follows pi; the name abbreviates \emph{reflective higher order},
and the transliteration is the pun. Quotation is written \quo{} and dropping
\drp{}; corner quotes are not used.

The design discussions that produced this calculus wrote the pair constructor
as an infix asterisk, $Q_p * Q_n$. That notation collides with the dropping
operator, so pairs are written here with angle brackets, $\dip{Q_p}{Q_n}$, and
the reader translating from those discussions should read $\dip{Q_p}{Q_n}$
wherever $Q_p * Q_n$ appears there.

% ---------------------------------------------------------------------------
\section{Why the namespace is expanded}
\label{sec:why}

The calculus below is the minimal response to four defects, each of which is a
defect of the \emph{unexpanded} namespace. They are recalled because each
feature of the design is answerable to one of them, and a reader who forgets the
defects will find the features arbitrary.

\paragraph{The scalars have no additive inverses.}
In the ordinary \rhoc\ calculus, structural congruence is associativity,
commutativity and a unit for parallel composition. Consequently $\Proc/\scong$
is the free commutative monoid on the non-parallel components of processes, and
$\Name$, being $\Proc/\scong$ transported through quotation, is that same free
commutative monoid. A free commutative monoid has no inverses and no
cancellation. Whatever else it can express, it cannot express the vanishing of
a sum of nonzero contributions. Since destructive interference is exactly such a
vanishing, no coefficient structure carried by names can produce it.

\paragraph{The multiplicative and additive units collide.}
Name multiplication is $\quo{}P \cdot \quo{}Q := \quo{}(P \mid Q)$, whose unit
is $\quo{}\pzero$. Any addition on names built from summation also has
$\quo{}\pzero$ for its unit. The rig law $0 \cdot x = 0$ then fails, since
$\quo{}\pzero \cdot x = x$. There is no absorbing element, and hence no rig.

\paragraph{There is no conjugation.}
Parallel composition is commutative, so any pairing built from it is symmetric
rather than Hermitian, and there is no involution available on the coefficients
to repair this. An involution must come from somewhere; the unexpanded calculus
has none of order two acting on scalars.

\paragraph{Multiplication by a monomial is a shift, not a scaling.}
This is the sharpest of the four. In the free commutative monoid, multiplying a
name by another name is injective and never decreases anything; it permutes
basis elements rather than combining them. Any attempt to realise the imaginary
unit as an ordinary name --- for instance by adjoining a ground process
$\mathsf{1}$ with $\mathsf{1}\mid\mathsf{1}\mid\mathsf{1}\mid\mathsf{1} \scong
\pzero$, so that $\quo{}\mathsf{1}$ has order four --- therefore fails for a
reason that is fatal rather than technical: $\quo{}R$ and
$\quo{}(\mathsf{1}\mid R)$ are names of \emph{different addresses}, so no
receiver hears both and nothing superposes. Multiplication by $i$ must move a
message \emph{between sectors at one address}, not between addresses. This is
the requirement that forces the sector to be a coefficient attached to a name
rather than a factor inside it, and it is the reason the design expands the
namespace instead of enriching the process language.

\begin{remark}
The four defects are not independent. Repairing the first by group completion
supplies the involution needed for the third, and separates the units demanded
by the second. Only the fourth is a genuinely separate requirement, and it is
the one that dictates \emph{where} the new structure is placed.
\end{remark}

% ---------------------------------------------------------------------------
\section{Syntax}
\label{sec:syntax}

\subsection{Sorts}

The calculus has three sorts: processes, pairs and names. Pairs are the new
sort. A pair is intended as a formal difference of processes, its first
component positive and its second negative. Pairs occur in two places only, on
messages and inside quotations; binders do not mention them, for the reason
given in Remark~\ref{rem:onebinder}.

\begin{definition}[Grammar]
\label{def:grammar}
\begin{align*}
  P, Q &::= \pzero \;\mid\; P \mid Q \;\mid\;
            \mathsf{for}(y \leftarrow x)P \;\mid\;
            x!A \;\mid\; \drp{l}x \;\mid\; \drp{r}x \\
  A, B &::= \dip{P}{Q} \\
  x, y &::= \quo{l}A \;\mid\; \quo{r}A
\end{align*}
\end{definition}

Three points of reading. A send carries a pair, not a process. A receipt binds a
single name variable, and the negative component is reached from it by
$\negop$, which Proposition~\ref{prop:lift} makes available on every term. And a
name is a quotation of a pair, tagged with a sector $l$ or $r$.

\begin{convention}[Drops in process position]
\label{con:drop}
$\drp{s}x$ is stipulated above to be a process, whereas the pair quoted by $x$
has two components. We take $\drp{s}x$ to denote the \emph{positive} component
of the $s$-sector content of $x$, so that $\drp{s}(\negop x)$ is the negative
component. No generality is lost and the ordinary reading of $\drp{}$ is
preserved.
\end{convention}

\begin{convention}[Abbreviations]
We write $x!Q$ for $x!\dip{Q}{\pzero}$. Together with
Convention~\ref{con:drop} and the choice $\quo{} := \quo{l}$, $\drp{} :=
\drp{l}$, this is what makes the ordinary \rhoc\ calculus a sub-language; see
Section~\ref{sec:conservativity}.
\end{convention}

\subsection{Negation}

\begin{definition}[Negation on pairs and names]
\label{def:neg}
$\negop$ is defined on pairs and on names by component exchange at the top:
\[
  \negop\dip{P}{Q} := \dip{Q}{P}, \qquad
  \negop\quo{s}A := \quo{s}(\negop A) \quad (s \in \{l,r\}).
\]
The definition does not recurse into the components.
\end{definition}

\begin{proposition}[Negation is a total involution on names]
\label{prop:involution}
$\negop$ is defined on every name without reference to reduction, and
$\negop\negop x = x$.
\end{proposition}

\begin{proof}
Immediate from component exchange being an involution on ordered pairs.
\end{proof}

\begin{proposition}[Unique lifting to processes]
\label{prop:lift}
SETTLED. There is exactly one extension of $\negop$ from names to processes
which is a homomorphism of the term algebra, namely
\[
  \negop\pzero := \pzero, \quad
  \negop(P \mid Q) := \negop P \mid \negop Q, \quad
  \negop(x!\dip{P}{Q}) := (\negop x)!\dip{\negop P}{\negop Q},
\]
\[
  \negop(\mathsf{for}(y \leftarrow x)P) := \mathsf{for}(y \leftarrow \negop x)\negop P,
  \quad
  \negop(\drp{s}x) := \drp{s}(\negop x).
\]
It satisfies $\negop\negop P = P$, and it descends to $\scong$ and to $\nameq$.
\end{proposition}

\begin{proof}
Processes are freely generated by the term formers over names, so a map on names
extends uniquely to an endomorphism of the term algebra; the displayed clauses
are that extension. Involutivity follows by induction from
Proposition~\ref{prop:involution}. For $\scong$, the clauses are homomorphic and
leave the binder alone, so associativity, commutativity and the unit are
preserved. For $\nameq$, applying $\negop$ to the quote-drop law
(Definition~\ref{def:nameq}) at $x$ yields the quote-drop law at $\negop x$,
since $\drp{s}(\negop x)$ is the negative component of $x$.
\end{proof}

\begin{remark}[What the lift is, and is not]
\label{rem:whatlift}
It is worth being explicit, because the notation invites a conflation that the
rest of the note depends on avoiding. The lifted $\negop$ is an
\emph{automorphism of the term algebra} --- a relabelling of the namespace ---
and not an additive inverse. In particular $P \mid \negop P$ is inert and is not
$\pzero$: a send and its mirror have no receipt between them, so nothing
reduces. Cancellation does not come from $\negop$ on processes; it comes from
the completion equations on pairs, and orthogonality accordingly comes from
where Remark~\ref{rem:normalisation} says it does.

Homomorphism for $\mid$ is therefore a \emph{consequence} of freeness, not a
design decision, and it carries none of the algebraic content it would carry if
$\negop$ were an inverse.
\end{remark}

\begin{remark}[Why the pair is inside the quotation]
\label{rem:inside}
SETTLED. If pairs occurred only on messages, then names would quote single
processes and $\negop x$ would be undefined: negation would be an operation on
messages that the namespace cannot express, so $x$ and $\negop x$ would not both
be addresses, the polarity of a participant in a communication could not be
stated, and Proposition~\ref{prop:lift} would have nothing to lift. Placing the
pair inside the quotation makes negation an operation on addresses and lets the
reduction rule mention $\negop x$ as a name like any other.

There is a second reason, of a different kind. Under the reading in which
parallel composition is addition on processes and name multiplication is
$\quo{}P \cdot \quo{}Q = \quo{}(P\mid Q)$, quotation is an exponential.
Completing to formal differences \emph{before} exponentiating gives a namespace
that is the exponential image of a group; completing afterwards does not, and
that is why earlier attempts to obtain negation dynamically, as a consequence of
annihilation, kept having to produce the sign at reduction time rather than
having it available in the syntax.
\end{remark}

% ---------------------------------------------------------------------------
\section{Sectors, polarity, and the group of a name}
\label{sec:group}

\subsection{Two quotations force two drops}

\begin{proposition}[Collapse]
\label{prop:collapse}
Suppose the calculus has two quotation operators $\quo{l}, \quo{r}$ and a
\emph{single} dropping operator $\drp{}$ satisfying $\drp{}(\quo{s}A) = A$ for
both $s$, together with the quote-drop law $\quo{s}(\drp{}x) \nameq x$. Then
$\quo{l}A \nameq \quo{r}A$ for every $A$.
\end{proposition}

\begin{proof}
Instantiate the quote-drop law at $x := \quo{r}A$ and $s := l$:
$\quo{l}A = \quo{l}(\drp{}(\quo{r}A)) \nameq \quo{r}A$.
\end{proof}

So the two quotations are distinguishable only if the drops are split as well;
a map has at most one two-sided inverse. This is why
Definition~\ref{def:grammar} carries $\drp{l}$ and $\drp{r}$ and not a single
$\drp{}$.

\begin{definition}[Sector projection]
\label{def:proj}
$\drp{s}(\quo{t}A) := A$ if $s = t$, and $\dip{\pzero}{\pzero}$ if $s \neq t$.
\end{definition}

\begin{remark}
FORK. Definition~\ref{def:proj} makes the drops total, so that $\drp{l}$ and
$\drp{r}$ behave as real and imaginary part. The alternative --- leaving the
cross cases stuck --- makes the drops partial and puts the sector distinction
outside the equational theory, where the matcher cannot use it. Totality is
taken here, but the choice interacts with Obligation~\ref{ob:observability}
below and is not innocent.
\end{remark}

\subsection{The four names of an address}

Each name has an underlying pair and two binary attributes: its sector, $l$ or
$r$, and its polarity, which is the choice between a pair and its exchange.
Fixing the underlying pair up to exchange therefore leaves four names.

\begin{definition}[Address and phase]
\label{def:addr}
For $x = \quo{s}A$, write $\sect{x} := s$ and let $\addr{x}$ be the pair $A$
taken up to $\negop$. The four names sharing an address are
\[
  \quo{l}A, \qquad \quo{r}A, \qquad \quo{l}(\negop A), \qquad \quo{r}(\negop A).
\]
\end{definition}

\begin{definition}[The imaginary unit]
\label{def:i}
SETTLED, with the caveat of Remark~\ref{rem:klein}. Multiplication by $i$ is the
successor map on the cycle
\[
  \quo{l}A \;\longmapsto\; \quo{r}A \;\longmapsto\; \quo{l}(\negop A)
  \;\longmapsto\; \quo{r}(\negop A) \;\longmapsto\; \quo{l}A .
\]
Equivalently: passing from $l$ to $r$ leaves the pair alone, and passing from
$r$ to $l$ exchanges it.
\end{definition}

\begin{proposition}
\label{prop:z4}
The map of Definition~\ref{def:i} has order four, its square is $\negop$, and the
four names of an address form a torsor under $\mathbb{Z}/4$ with $\phase{-}$ the
resulting coordinate, $\phase{\quo{l}A} = 1$.
\end{proposition}

\begin{proof}
Two applications carry $\quo{l}A$ to $\quo{l}(\negop A)$, which is $\negop$ by
Definition~\ref{def:neg}; $\negop$ is an involution by
Proposition~\ref{prop:involution}, so the map has order four. The action is free
and transitive on the four names by Definition~\ref{def:addr}.
\end{proof}

\begin{remark}[Why $i$ is not the composite of two independent involutions]
\label{rem:klein}
This is the most delicate point in the design and it must not be glossed.
Sector exchange and pair exchange are each involutions, and if they are taken as
\emph{independent} the group they generate is the Klein four-group. The Klein
four-group is the group of the split-complex numbers, in which $j^2 = +1$ and
there are zero divisors, $(1+j)(1-j) = 0$. Zero divisors are fatal here, since a
pairing would then vanish for reasons having nothing to do with orthogonality.
Definition~\ref{def:i} is precisely the stipulation that the two exchanges are
\emph{not} independent: the descending passage from $r$ to $l$ carries the pair
exchange with it, so that the generated group is cyclic of order four rather
than Klein. The asymmetry is not decoration. It is the content of the
requirement $i^2 = -1$, and it is what a design with two free involutions cannot
supply.
\end{remark}

\begin{remark}[Only one doubling]
Name multiplication is parallel composition, hence commutative. The
Cayley--Dickson construction preserves commutativity once and loses it at the
second step. A namespace whose product is parallel composition therefore admits
exactly one doubling and no more: the calculus can be complex, and cannot be
quaternionic, for a structural reason rather than by stipulation.
\end{remark}

% ---------------------------------------------------------------------------
\section{Structural theory}
\label{sec:structural}

\begin{definition}[Structural congruence]
\label{def:scong}
$\scong$ is the least congruence containing alpha-equivalence and satisfying
associativity, commutativity and $\pzero$ as unit for $\mid$, together with
\[
  \negop\negop A \scong A .
\]
\end{definition}

\begin{definition}[Name equivalence]
\label{def:nameq}
$\nameq$ is the least congruence satisfying
\[
  \quo{s}\dip{\drp{s}x}{\negop\drp{s}x} \nameq x \quad
    \text{when } \sect{x} = s,
  \qquad
  \frac{A \scong B}{\quo{s}A \nameq \quo{s}B}.
\]
\end{definition}

\begin{remark}[The completion equations, and where they live]
\label{rem:completion}
FORK, and the most consequential one in the note. A pair is intended as a formal
difference, which suggests two further equations:
\[
  \dip{Q}{Q} \scong \pzero,
  \qquad
  \dip{P}{Q} \scong \dip{P \mid R}{Q \mid R}.
\]
These are what make the pair sort the group completion of processes rather than
merely an ordered pair. But $\nameq$ contains $\scong$ under quotation, and COMM
fires on $\nameq$. Imposing them therefore makes them equations \emph{on
addresses}, which the matcher must decide.

The two readings differ sharply.
\begin{itemize}
\item \textbf{Arithmetic names.} The completion equations hold on names.
  Cancellation is then real, $\negop$ is genuine negation, and destructive
  interference is the matcher discovering that a pair's components have become
  congruent. The cost is that unique factorisation is lost, the monomial basis
  with it, and the decidability of name matching becomes a serious question.
\item \textbf{Coloured names.} The completion equations do not hold on names,
  which are then a free set with a $\mathbb{Z}/4$ action. Negation is a label
  carried along and read by the reduction rule, exactly as the sector is. The
  matcher is untouched; the arithmetic is deferred to a coefficient layer
  outside the calculus.
\end{itemize}
The reduction rule of Section~\ref{sec:comm} never asks whether two
differently-signed names are equal, only whether they share an address, so the
coloured reading suffices for the operational semantics as it stands. The
arithmetic reading is what the semantic programme wants. The note takes neither;
it records that the choice determines whether cancellation is a structural event
or an external one.
\end{remark}

% ---------------------------------------------------------------------------
\section{Reduction}
\label{sec:comm}

\subsection{Matching is sector-blind and polarity-blind}

\begin{definition}[Address matching]
\label{def:match}
SETTLED. A receipt at $x$ and a send at $x'$ are \emph{matched} when
$\addr{x} = \addr{x'}$, irrespective of $\sect{x}, \sect{x'}$ and irrespective
of polarity.
\end{definition}

This is forced by the fourth defect of Section~\ref{sec:why}. If $\quo{l}A$ and
$\quo{r}A$ were distinct addresses, then $\quo{l}A!Q \mid \quo{r}A!Q$ would be
two messages that no single receiver can hear, rather than a superposition of
one message with itself in two sectors; nothing would ever combine, the sector
would be a superselection rule, and the doubling would carry no information.
Sector and polarity are coefficients on a name, not part of its identity.

\subsection{The rule}

\begin{definition}[Relative polarity]
\label{def:relpol}
For matched $x$ and $x'$, the \emph{relative polarity} $\delta(x,x')$ is the
identity when $x' = x$ and $\negop$ when $x' = \negop x$. It is well defined by
Definition~\ref{def:addr}, since matched names differ by an element of the
$\mathbb{Z}/4$ torsor and polarity is the component of that element which
Definition~\ref{def:i} squares to.
\end{definition}

\begin{definition}[COMM]
\label{def:comm}
Let $t := \sect{x}$ and let $\mathsf{op}(l) := r$, $\mathsf{op}(r) := l$. For
matched $x, x'$ in the sense of Definition~\ref{def:match},
\[
  \mathsf{for}(y \leftarrow x)P \;\mid\; x'!A
  \;\red\;
  P\{\, \quo{\mathsf{op}(t)}\bigl(\delta(x,x')\,A\bigr) \,/\, y \,\}
\]
The rule is closed under parallel composition and under $\scong$ in the usual
way.
\end{definition}

One clause, two effects, and they remain independent of one another.

\paragraph{Sector: communication intertwines.}
SETTLED. The quotation deposited is in the sector \emph{opposite} to the one the
receipt listened on. A communication on the real axis deposits an imaginary
name; a communication on the imaginary axis deposits a real one. So a
communication event is the map $l \mapsto r \mapsto l$, and by
Definition~\ref{def:i} it is multiplication by $i$.

\paragraph{Polarity: opposition exchanges the pair.}
SETTLED. When the participants have opposite polarity, $\delta$ is $\negop$ and
the deposited pair is exchanged, which by Definition~\ref{def:neg} is
multiplication by $\negop$. So opposite polarity contributes the sign, and
$i^2 = \negop 1$ is realised as: two sector rotations equal one polarity
exchange.

\begin{remark}[One binder, one clause]
\label{rem:onebinder}
SETTLED, and the reason the grammar of Definition~\ref{def:grammar} binds a
single variable. An earlier presentation of this calculus had receipts bind a
pair $\dip{y_p}{y_n}$ and stated COMM in two clauses, one per relative polarity,
each depositing $\quo{}\dip{Q_p}{\pzero}$ and $\quo{}\dip{Q_n}{\pzero}$ into the
two variables. Both features are removed by Proposition~\ref{prop:lift}. The
whole received pair goes to the single variable $y$, and the negative component
is $\drp{s}(\negop y)$, so the second binder was never carrying information the
first did not. The polarity clause then ceases to be a second case and becomes
the factor $\delta$, which is the statement that the relative group element of
the two participants is the phase deposited --- visible in the rule rather than
distributed over two of them.
\end{remark}

\begin{proposition}[Equivariance]
\label{prop:equivariance}
$P \red P'$ implies $\negop P \red \negop P'$.
\end{proposition}

\begin{proof}
$\negop$ is a homomorphism leaving binders alone
(Proposition~\ref{prop:lift}), so it carries a redex
$\mathsf{for}(y \leftarrow x)P \mid x'!A$ to
$\mathsf{for}(y \leftarrow \negop x)\negop P \mid (\negop x')!(\negop A)$, which
is a redex: matching is polarity-blind (Definition~\ref{def:match}) and
$\delta(\negop x, \negop x') = \delta(x,x')$. Both reducts are
$\negop P\{\quo{\mathsf{op}(t)}(\delta(x,x')\,\negop A)/y\}$, using
$\negop\quo{s}B = \quo{s}(\negop B)$ and the commutation of $\negop$ with
substitution, which holds because $\negop$ acts uniformly on the names of $P$.
\end{proof}

\begin{remark}[Equivariance is a selection principle]
\label{rem:selection}
Proposition~\ref{prop:equivariance} is not decoration; it is what picks the
deposit convention out of the alternatives. Requiring reduction to commute with
the unique lifting of Proposition~\ref{prop:lift} rules out the earlier
convention, under which the deposited name was always
$\quo{\mathsf{op}(t)}\dip{Q}{\pzero}$ with the payload in the positive slot:
nothing of that shape lies in the image of an equivariant rule, since
$\negop\quo{}\dip{Q}{\pzero} = \quo{}\dip{\pzero}{Q}$, and the failure is
uniform at one application of $\negop$ per communication. Depositing the whole
pair is the unique repair, and it is the convention of
Definition~\ref{def:comm}. That the same convention also collapses the pair
binder and the second COMM clause is the reason to prefer it on grounds other
than symmetry.
\end{remark}

\begin{remark}[Phase is path length]
\label{rem:pathlength}
An immediate consequence of the sector effect. The sector of a name is the
parity of the number of communications that produced it, so the relative phase
between two derivations reaching a common outcome is the parity of the
difference of their lengths. This is worth naming, because the question of where
destructive interference could come from in an unmodified \rhoc\ calculus had
the answer \emph{length-asymmetric reconvergence}, with no mechanism to supply
it: concurrency diamonds are symmetric, so the holonomy vanishes. Here the phase
\emph{is} the length asymmetry, with no cost functor, no weight map and no
external coefficient algebra.
\end{remark}

\begin{remark}[The identity step is not fictitious]
Under the reading of Remark~\ref{rem:pathlength}, a step that changes no process
still rotates the sector. A device introduced elsewhere to complete deficient
rows, and justified there case by case, acquires an independent justification
here: doing nothing and doing nothing once are genuinely different.
\end{remark}

\subsection{The four-case refinement}
\label{sec:fourcase}

\begin{remark}
FORK. Definition~\ref{def:comm} computes the deposited sector from the
receipt's channel alone. But Definition~\ref{def:match} makes matching
sector-blind, so a receipt at $\quo{l}A$ may be served by a send at $\quo{r}A$,
and the rule as stated ignores the sender's sector. The natural completion is
that the two sectors \emph{multiply}:
\[
\begin{array}{c|cc}
  & \text{send } l & \text{send } r \\ \hline
  \text{recv } l & 1 & i \\
  \text{recv } r & i & -1
\end{array}
\]
with $-1$ realised as $\negop$ on the deposited pair and the off-diagonal entries
being exactly the intertwining of Definition~\ref{def:comm}. Under this reading
the rule of Definition~\ref{def:comm} is the off-diagonal half of the table, and
$i^2 = -1$ is a case of the reduction rule rather than a separate stipulation.

The refinement is attractive and is not adopted here, because it changes what
the rule computes rather than extending it: on the diagonal it deposits a name
in the sector it received on, which contradicts the plain reading of
``communication intertwines''. Settling this is the first item of
Section~\ref{sec:obligations}.
\end{remark}

\begin{remark}[Amplitude is multiplicity]
Under Definition~\ref{def:match}, the expression
$\quo{l}A!Q \mid \quo{r}A!Q$ is a single address carrying one copy in each
sector, and with the pair supplying signs this is $(1+i)$ times a send at $A$.
Generally, the coefficient of a message is the multiset of sectors and
polarities its copies occupy, so Gaussian-integer coefficients are realised as
multiplicity under parallel composition, with no external coefficient algebra
and no tagging. This is the payoff that motivated
Definition~\ref{def:match}, and it should be the first thing checked against any
proposed revision of the rule.
\end{remark}

% ---------------------------------------------------------------------------
\section{Conservativity}
\label{sec:conservativity}

\begin{proposition}[Embedding]
\label{prop:conservativity}
FORK, stated as a target rather than a theorem. Let
$\iota(-)$ send an ordinary \rhoc\ term to the term of
Definition~\ref{def:grammar} obtained by
\begin{align*}
  \iota(x!Q) &:= \iota(x) ! \dip{\iota(Q)}{\pzero},
  &
  \iota(\quo{}P) &:= \quo{l}\dip{\iota(P)}{\pzero}, \\
  \iota(\mathsf{for}(y \leftarrow x)P) &:=
    \mathsf{for}(y \leftarrow \iota(x))\iota(P),
  &
  \iota(\drp{}x) &:= \drp{l}\iota(x)
\end{align*}
homomorphically elsewhere. Then $\iota(-)$ is
injective and preserves $\scong$.
\end{proposition}

\begin{remark}[Conservativity fails for reduction, and deliberately]
\label{rem:notconservative}
The image of $\iota(-)$ is not closed under $\red$. A single
communication moves a name off the real axis, so an ordinary \rhoc\ program is
not, in this calculus, a program that stays classical. This is not an oversight
to be repaired: it is the content of Remark~\ref{rem:pathlength}, since a
calculus in which ordinary programs remained on the real axis would be one in
which no phase ever accumulated. The design is a genuinely different calculus,
not an extension, and it should be presented as such.

If a conservative reduction is wanted, the diagonal entries of the table in
Section~\ref{sec:fourcase} supply it: under that reading, a real receipt served
by a real send stays real, and the ordinary calculus embeds as the real
sub-calculus. That is an argument for the refinement, and is recorded here
rather than in Section~\ref{sec:fourcase} because it is the strongest one.
\end{remark}

% ---------------------------------------------------------------------------
\section{Evaluation and residuation}
\label{sec:odot}

The calculus above is a term language and a reduction relation. The construction
it serves needs two further operations: an evaluation, which runs a composite to
its terminal states and collects them, and a pole, against which the dagger is
obtained by residuation. Both are recorded here. Following the design decision
that the pole be taken as simply as possible for the time being, $\bot$ is the
zero process and $P^{\perp}$ is residuation to it --- the linear perp.

\subsection{Inert states and evaluation}

\begin{definition}[Inert]
\label{def:inert}
$P$ is \emph{inert} when $P \not\red$; that is, when no COMM applies to any of
its components under $\scong$. Write $\mathsf{Inert}$ for the inert processes.
\end{definition}

Inertness is not triviality: $\pzero$ is inert, and so is $x!A$ with no matching
receipt. The distinction matters throughout what follows, and is the reason
$\bot$ below is not simply ``is inert''.

\begin{definition}[Evaluation product]
\label{def:odot}
\[
  P \odot Q \;:=\; \sum_{\substack{R \in \mathsf{Inert} \\ P \mid Q \;\red^{*}\; R}} R
\]
where the sum ranges over terminal states with multiplicity, one summand per
\emph{causal history} ending at $R$.
\end{definition}

\begin{proposition}[Commutativity]
\label{prop:odotcomm}
$P \odot Q \scong Q \odot P$.
\end{proposition}

\begin{proof}
$\mid$ is commutative and the index of the sum mentions $P$ and $Q$ only through
$P \mid Q$.
\end{proof}

\begin{definition}[Evaluation]
\label{def:eval}
$\Downarrow P := \pzero \odot P$.
\end{definition}

\begin{remark}[Histories, not paths]
\label{rem:histories}
SETTLED, and load-bearing rather than hygienic. Indexing the sum by reduction
\emph{sequences} rather than by causal histories counts the linearisations of a
partial order, so $m$ independent redexes contribute a spurious factor of $m!$.
The multiplicity in Definition~\ref{def:odot} is therefore the number of
histories, equivalently the number of maximal chains in the reduction order
modulo permutation of independent steps. Concurrency squares are filled; what is
summed is what survives the quotient.
\end{remark}

\begin{remark}[What $\odot$ is not]
\label{rem:odotalgebra}
Three properties fail, and each failure is informative rather than repairable.

\emph{No unit.} $\pzero \odot P = \Downarrow P$, which is $P$ only when $P$ is
already inert. So $\pzero$ is a unit on $\mathsf{Inert}$ and nowhere else, which
is one reason the carrier of the eventual algebra must be cut down --- see
Section~\ref{sec:pole}.

\emph{Not idempotent.} $P \odot P \neq P$ in general, because two copies of a
process interact. This is a feature: it is what makes multiplicity meaningful,
and it is why the additive structure cannot be mixed choice, whose idempotence
would force $2 = 1$.

\emph{Not associative.} $(P \odot Q) \odot S$ forces $P \mid Q$ to reach a
terminal state before $S$ arrives, so any global reduction of $P \mid Q \mid S$
that schedules a communication between $P$ and $S$ early, destroying a
$P$--$Q$ redex, is a history of the unbracketed composite and of no bracketed
one. The bracketed product undercounts, and the deficit is exactly the
cross-redexes. Writing $M_R$ for the operator $\square \mid R$ and $T$ for
one-step reduction, the deficit is the commutator $[T, M_R]$: the redexes
between $R$ and the state that existed in neither alone. $M_R$ alone is diagonal
in the occupation basis and commutes with every $M_S$; $T$ alone is the
off-diagonal part; neither is interesting, and their failure to commute is where
the content of the construction lives.
\end{remark}

\begin{convention}[The $n$-ary form]
\label{con:nary}
Because of the last item, the primitive is not the binary product but the
family of symmetric $n$-ary evaluations
\[
  \odot(R_1, \ldots, R_n) \;:=\;
  \sum_{\substack{R \in \mathsf{Inert} \\ R_1 \mid \cdots \mid R_n \;\red^{*}\; R}} R
\]
with one global evaluation and no bracketing. Symmetry under permutation is
Proposition~\ref{prop:odotcomm}. Nothing needs to be associative because nothing
is composed twice; factorisation of an $n$-ary evaluation into lower ones then
becomes a property that holds exactly on causally independent components, rather
than an axiom.
\end{convention}

\begin{remark}[What the sum is]
\label{rem:sumfork}
FORK. Under the reading in which parallel composition is addition, the formal
sum of Definition~\ref{def:odot} is available in the calculus: it is the
parallel composition of the terminal states. That reading is attractive, since
it needs no new former and no tagging, and multiplicity is then counted by
$\mid$ itself.

It has a hazard that must be checked before it is adopted. Inertness is not
preserved by parallel composition: $x!A$ and
$\mathsf{for}(y \leftarrow x)P$ are each inert and their
composition is not. So the parallel composition of the summands may reduce
further, and the sum would not be a normal form. The alternative is an external
formal sum over the multiset of terminal states, which is inert by fiat and
costs a coefficient layer outside the calculus.

Two things are settled regardless. Tagging the summands with fresh names is
\emph{not} the way to defeat idempotence: a retained tag is a which-path record,
distinguishing summands permanently and preventing exactly the coincidence of
outcomes on which the addition of coefficients depends. And quotation carries
$\mid$ to the name product, so quoting an internalised sum yields a
\emph{product} of quotes; whichever reading is taken, the sum must not be
internalised beneath a quotation.
\end{remark}

\begin{remark}[Partiality]
\label{rem:partiality}
$\odot$ is partial and the sum may be infinite. Non-confluence is not itself an
obstacle --- the sum is over all terminal states, not over a chosen one --- but
divergence is: a composite with no terminal state contributes the empty sum, and
a composite with infinitely many contributes a sum that requires a summability
condition before it denotes anything. Both are recorded as
Obligation~\ref{ob:summability}.
\end{remark}

\subsection{The pole, and residuation}
\label{sec:pole}

\begin{definition}[Residual]
\label{def:residual}
For formulae $\varphi, \psi$ of the generated spatial logic, $\varphi \rhd \psi$
is the right adjoint to the separating conjunction:
\[
  P \models \varphi \rhd \psi
  \quad\text{iff}\quad
  \forall R.\; R \models \varphi \;\Rightarrow\; P \mid R \models \psi .
\]
\end{definition}

\begin{definition}[The pole]
\label{def:pole}
$\bot$ is the behavioural formula
\[
  P \models \bot \quad\text{iff}\quad \Downarrow P \scong \pzero,
\]
that is, every terminal state of $P$ is the zero process.
\end{definition}

\begin{definition}[Linear perp]
\label{def:perp}
$P^{\perp} := \{\, Q \;:\; P \odot Q \scong \pzero \,\}$, and for a set
$S$ of processes, $S^{\perp} := \bigcap_{P \in S} P^{\perp}$.
\end{definition}

\begin{proposition}[$P^{\perp}$ is residuation to $\bot$]
\label{prop:perpisres}
$Q \in P^{\perp}$ iff $Q \models \psi_P \rhd \bot$, where $\psi_P$ is a formula
whose extension is the equivalence class of $P$ in the logic.
\end{proposition}

\begin{proof}
By Definition~\ref{def:residual}, $Q \models \psi_P \rhd \bot$ iff
$Q \mid R \models \bot$ for every $R$ in the extension of $\psi_P$; by
Definition~\ref{def:pole} that is $\Downarrow(Q \mid R) \scong \pzero$, which is
$R \odot Q \scong \pzero$ by Definitions~\ref{def:odot} and~\ref{def:eval}.
\end{proof}

\begin{corollary}[Symmetry]
$Q \in P^{\perp}$ iff $P \in Q^{\perp}$.
\end{corollary}

\begin{definition}[States]
\label{def:states}
$P$ is a \emph{state} when $\{P\}^{\perp\perp} = \{P\}^{\perp\perp}$ is
generated by $P$, that is, when $P$ is $\perp\perp$-closed. The states are the
carrier of everything that follows.
\end{definition}

\begin{remark}[The closure is the null-space quotient]
\label{rem:closure}
$\{P\}^{\perp\perp}$ is the set of processes annihilated by every annihilator of
$P$, so biorthogonal closure is closure under the testing preorder induced by
the pole, with the tests being annihilators and success being reduction to the
zero process. Three consequences.

It supplies the collapse that quotation cannot. Quotation is injective on
$\Proc/\scong$ and therefore identifies nothing; every identification in the
eventual pairing comes from $\odot$ and from this closure, that is, from
reduction.

It fixes the equivalence at which the construction lives, and the equivalence it
fixes is \emph{testing}, not bisimulation. Since $\nameq$ is intensional and
testing is not, the coefficient object and the state space sit at different
equivalences; that is a known tension and it is not resolved here.

It destroys freeness. $\Proc/\scong$ under $\mid$ is free commutative on the
non-parallel components, and closed sets are not freely generated. The
occupation basis is a coordinate system for the unclosed calculus and does not
survive into the closed one.
\end{remark}

\begin{remark}[Totality fails, and how it is repaired]
\label{rem:totality}
Not every process has an annihilator: a persistent or replicating process at a
name no one else holds has none, so $\psi_P \rhd \bot$ is unsatisfiable and the
dagger is partial. Definition~\ref{def:states} is the repair adopted here ---
restrict to the closed elements, which is the standard move --- and it is the
reason the states are cut down rather than the pole loosened. The alternative,
kept in reserve, is to relativise the pole to a residue, taking $R \vdash P
\perp Q$ when every terminal state of $P \mid Q$ is $R$; that would make the
pairing residue-valued in exactly the way the intended semantics wants, at the
cost of a family of poles rather than one.
\end{remark}

\begin{remark}[Canonical witnesses are forced]
\label{rem:canonical}
$P^{\perp}$ is a set, and the dagger must be a term, because it is going to be
quoted. Quotation does not respect any behavioural equivalence, so a
representative must be chosen and not merely existentially quantified. The
choice adopted is by minimal structural complexity, the measure already present
for names and processes, with ties broken by the summation of
Definition~\ref{def:odot} when the minimum is not unique.
\end{remark}

\subsection{The pairing}

\begin{definition}[Dagger and pairing]
\label{def:pairing}
For a state $P$, $P^{\dagger}$ is the canonical witness of $\psi_P \rhd \bot$ in
the sense of Remark~\ref{rem:canonical}, composed with the sector and polarity
correction of Section~\ref{sec:dagger}. The pairing is
\[
  \langle P \mid Q \rangle \;:=\; \quo{l}\dip{P^{\dagger} \odot Q}{\pzero}
\]
and, for a reified operator, using Convention~\ref{con:nary},
\[
  \langle P \mid K \mid Q \rangle \;:=\;
  \quo{l}\dip{\odot\!\left(P^{\dagger},\, K^{*}_{x},\, x!\dip{Q}{\pzero}\right)}{\pzero}
\]
with $x$ private to the composite and $K^{*}_{x} :=
\mathsf{for}(y \leftarrow x)K[\drp{l}y]$ the abstraction
corresponding to $K$.
\end{definition}

\begin{remark}[Why the quotation is outermost, and why the operator is reified]
\label{rem:outermost}
Two points, each answering an earlier error.

$\odot$ lands in $\Proc$ and the scalars are names, so the pairing must pass
through a quotation to change sort; a definition that has $\odot$ produce a
scalar directly collapses the two sorts at precisely the place where the
collapse is dangerous. The quotation is applied once, after evaluation, and
contributes no identifications of its own.

Reifying $K$ removes the bracketing question rather than answering it. Filling a
hole and then evaluating, or evaluating and then filling, differ by the
cross-redexes of Remark~\ref{rem:odotalgebra}; passing the argument along a
private wire means there is no hole-filling anywhere, and the bracket is
manifestly the three-point evaluation it should have been. Privacy of $x$ is
required, not cosmetic: it is what prevents $P^{\dagger}$ from intercepting the
operator's wire, and renaming cannot supply it in a calculus where any party may
construct any quotation.
\end{remark}

\begin{remark}[Annihilation is normalisation, not orthogonality]
\label{rem:normalisation}
An immediate and initially unwelcome consequence. $\langle P \mid P \rangle =
\quo{l}\dip{\pzero}{\pzero}$, the unit of the name product, so annihilation
gives $1$. The pole is a \emph{normalisation} condition: $P^{\dagger}$ is the
process that reduces $P$ to the unit, and $\langle P \mid P \rangle = 1$ is a
theorem rather than a stipulation. That is a gain, but it vacates the position
orthogonality was expected to occupy, and the correspondence table of the
earlier work, which reads ``orthogonality is process annihilation'', is
backwards on this point.

Orthogonality must therefore come from cancellation among the summands of
Definition~\ref{def:odot}, which requires additive inverses on those summands.
This is precisely what the pair sort supplies and what the unexpanded namespace
of Section~\ref{sec:why} could not: two histories reaching the same terminal
state with opposite polarity cancel, and $\langle P \mid Q \rangle = 0$ is the
vanishing of a sum of nonzero contributions rather than the absence of any. The
expansion of the namespace and the definition of the pairing are answerable to
each other at exactly this point.
\end{remark}

% ---------------------------------------------------------------------------
\section{The dagger, in one paragraph}
\label{sec:dagger}

The construction this calculus serves requires an antilinear involution on
states. Section~\ref{sec:odot} supplies the pole and the residuation; what
remains is the sector and polarity content, recorded here.

First, the dagger is not a syntactic polarity swap on term formers. It is
obtained by residuation against a pole: $Q^\dagger$ is a witness of
$\psi_Q \rhd \bot$, in the sense made precise by
Proposition~\ref{prop:perpisres} and Remark~\ref{rem:canonical}. This makes the
dagger a property first and a term second, hence stable under any presentation
generating the same logic. The choice of $\bot$ as the zero process is the
simplest available and is provisional; Remark~\ref{rem:totality} records the
relativised alternative.

Second, the sector and polarity components of the dagger are now fixed. Flipping
both sectors is the map $(a,b) \mapsto (b,a)$, which in the complex numbers is
$z \mapsto i \cdot \bar{z}$ and not conjugation; taken as the dagger it makes
the pairing skew-Hermitian and puts $\langle P \mid P \rangle$ in the imaginary
sector, so positivity fails outright. The correction is to compose the sector
flip with a pair exchange,
\[
  \quo{l} \mapsto \negop\quo{r}, \qquad \quo{r} \mapsto \negop\quo{l},
\]
which is $z \mapsto -i\bar{z}$, is involutive, and is assembled entirely from
operations Definitions~\ref{def:neg} and~\ref{def:i} already provide. Since
Proposition~\ref{prop:lift} lifts $\negop$ to processes, this composite acts on
a whole process rather than name by name, which is what the pairing of
Definition~\ref{def:pairing} requires and what an earlier presentation could not
write down. Remark~\ref{rem:whatlift} remains in force: $\negop$ supplies the
dagger's polarity component and is not itself the dagger.

\begin{remark}[The positivity obligation, made precise]
\label{rem:positivity}
Since every communication rotates the sector and the dagger rotates it once
more, the sector in which $\langle P \mid P \rangle$ lands depends on the number
of communications consumed in the annihilation of $P^\dagger$ against $P$,
modulo four. Positivity therefore requires that all annihilation paths of
$P^\dagger \mid P$ agree modulo four. This is a graded confluence condition,
weaker than confluence, and it is not obviously true. It is cheap to test on the
ladder of mutually annihilating names, and it should be tested before anything
further is built on the dagger.
\end{remark}

% ---------------------------------------------------------------------------
\section{Obligations}
\label{sec:obligations}

\begin{obligation}[Cross-redexes and the unit of evaluation]
\label{ob:assoc}
Remark~\ref{rem:odotalgebra} makes $\odot$ non-associative, with the deficit
$[T, M_R]$. Convention~\ref{con:nary} sidesteps the question for a single
evaluation but does not answer it: composing operators, $\langle P \mid K_1 K_2
\mid Q \rangle$, reintroduces it in the form of whether the wires create
cross-redexes. Determine the condition under which the $n$-ary evaluation
factors, and check it against the reified operators of
Definition~\ref{def:pairing}.
\end{obligation}

\begin{obligation}[Summability]
\label{ob:summability}
Discharge Remark~\ref{rem:partiality}: give the condition under which the sum of
Definition~\ref{def:odot} denotes when the set of terminal states is infinite.
Over a positive codomain monotone convergence suffices; with the pair sort
supplying signs it does not, and an absolute-summability condition on the
annihilation profile is the candidate.
\end{obligation}

\begin{obligation}[The sum, internal or external]
\label{ob:sum}
Settle Remark~\ref{rem:sumfork}. If the sum is parallel composition, show that
the composite of the terminal states is inert, or characterise when it is not.
If external, say where the coefficient layer lives and how it interacts with
Remark~\ref{rem:completion}.
\end{obligation}

\begin{obligation}[Matching the logic to the pole]
\label{ob:logicmatch}
Proposition~\ref{prop:perpisres} uses a formula $\psi_P$ whose extension is the
equivalence class of $P$. If that logic separates more finely than
$\perp\perp$-equivalence, then $P \mapsto P^{\dagger}$ does not factor through
the states of Definition~\ref{def:states} and the dagger is ill defined on them.
Either take the characteristic formula in a logic whose adequacy target is the
pole's equivalence, or define the dagger on terms and prove separately that it
descends.
\end{obligation}

\begin{obligation}[The four-case table]
Settle whether the deposited sector is computed from the receipt's channel alone
(Definition~\ref{def:comm}) or from both participants
(Section~\ref{sec:fourcase}). Remark~\ref{rem:notconservative} is the strongest
argument for the latter. This decision governs everything downstream and should
be taken first.
\end{obligation}

\begin{obligation}[Arithmetic or coloured names]
Settle Remark~\ref{rem:completion}: whether the completion equations hold on
names, and hence whether cancellation is a structural event decided by the
matcher or an external one deferred to a coefficient layer.
\end{obligation}

\begin{obligation}[Decidability of name matching]
\label{ob:decidability}
COMM fires on $\nameq$. Under the arithmetic reading of
Remark~\ref{rem:completion}, the matcher must normalise formal differences
nested through quotation. Termination is not obvious, and this is the practical
gate on the entire design.
\end{obligation}

\begin{obligation}[Freshness]
\label{ob:freshness}
The property that a quotation of $P$ is not free in $P$ is what obviates a
freshness operator and underwrites freshness by quotation. Its proof proceeds by
structural complexity. Both quotation operators must increase that measure,
which is routine; but the quote-drop law of Definition~\ref{def:nameq} relates a
name to a term built from its own drops, and the completion equations, if
adopted, shorten terms non-structurally. The mutual recursion of $\scong$ and
$\nameq$ was already the delicate part of the metatheory. That argument must be
redone rather than adapted.
\end{obligation}

\begin{obligation}[Observability of the sector]
\label{ob:observability}
$\drp{l}$ and $\drp{r}$ are term formers, so any context may apply both to a
name and see which yields $\pzero$. The sector is therefore classically
readable, which makes it a superselection rule rather than a phase, and makes
global phase detectable by a two-line context. The natural remedy is to admit
the drops in guard position only, where evaluation is pure and non-consuming,
and to let term-position dropping be $\drp{l}$ by convention. This needs
deciding, and it interacts with Definition~\ref{def:proj}.
\end{obligation}

\begin{obligation}[Scalar extraction and occurrence counts]
\label{ob:linearity}
For a race to interfere, its candidates must reach the same outcome with
different coefficients. Consuming from $\quo{l}A$ rather than $\quo{r}A$ yields
$P\{\quo{r}Q/y\}$ rather than $P\{\quo{l}Q/y\}$, and these are the same outcome
up to phase only if the sector can be extracted from the name to the term, that
is, only if there is a homomorphism from processes to $\mathbb{Z}/4$ under which
substitution multiplies. The natural candidate is sector degree, the product of
sectors over name occurrences --- but then the phase is raised to the power of
the number of occurrences of $y$, so extraction holds only when $y$ occurs
exactly once. Occurrence count is the exponent on the phase. Either linear use
of names bound from graded channels is required, or multiplicative phase in
occurrences is accepted and scalar extraction abandoned. This is a linearity
requirement of a different kind from the ones previously considered and declined,
and the earlier reasons for declining do not apply to it.
\end{obligation}

\begin{obligation}[Positivity]
Discharge or refute Remark~\ref{rem:positivity}: that all annihilation paths of
$P^\dagger \mid P$ agree in length modulo four.
\end{obligation}

\begin{obligation}[Two gradings on names]
This calculus grades names by sector. A separate line of work grades them by
binding structure, unforgeable names being those introduced by a name-restriction
operator and classical ones being quotations. The two gradings are independent,
so there are four kinds of name, and their interaction is unstated. The
conjecture worth testing is that sectors should live only in the unforgeable
namespace, which would also confine Obligation~\ref{ob:observability}.
\end{obligation}

\begin{obligation}[Metatheory]
Nothing in this note is proved beyond Propositions~\ref{prop:involution},
\ref{prop:collapse} and~\ref{prop:z4}. Subject reduction for the sector grading,
whether bisimulation is a congruence in the presence of two quotations, and
whether the sector is visible to bisimulation at all, are all open. The last is
pressing: an inert distinction that $\scong$ can see and $\sim$ cannot is exactly
the configuration in which the spatial layer strictly refines behaviour, and the
sign structure would then be behaviourally invisible.
\end{obligation}

% ---------------------------------------------------------------------------
\section{What this note does not do}

It does not give the interpretation map, and it does not establish that the
resulting scalars form a rig --- only that the four obstructions to their doing
so have each been given a place to be repaired. The evaluation and the pole of
Section~\ref{sec:odot} are stated, not developed: nothing is proved about
$\odot$ beyond commutativity, the pole is the simplest one available and is
provisional, and the summability question of Obligation~\ref{ob:summability} is
open, so the pairing of Definition~\ref{def:pairing} is a definition and not yet
a well-defined operation. It does not treat the join, whose generalisation to
graded messages requires a sum over matchings with coefficients in a group ring,
and whose normalisation may or may not remain the right unit. And it takes no
position on whether the resulting calculus is worth the cost of
Obligation~\ref{ob:decidability}, which is the question a reader should keep in
view throughout.

% ---------------------------------------------------------------------------
\begin{thebibliography}{9}

\bibitem{MR05} L.G. Meredith and M. Radestock.
  \emph{A reflective higher-order calculus.}
  Electronic Notes in Theoretical Computer Science, 141(5), 2005.

\bibitem{Milner} R. Milner.
  \emph{Communicating and Mobile Systems: the Pi-Calculus.}
  Cambridge University Press, 1999.

\bibitem{CairesCardelli} L. Caires and L. Cardelli.
  \emph{A spatial logic for concurrency.}
  Information and Computation, 186(2), 2003.

\bibitem{Caires04} L. Caires.
  \emph{Behavioral and spatial observations in a logic for the pi-calculus.}
  In Foundations of Software Science and Computation Structures, 2004.

\bibitem{Girard} J.-Y. Girard.
  \emph{Linear logic.}
  Theoretical Computer Science, 50(1), 1987.

\bibitem{AbramskyCoecke} S. Abramsky and B. Coecke.
  \emph{A categorical semantics of quantum protocols.}
  In Logic in Computer Science, 2004.

\bibitem{CoeckeKissinger} B. Coecke and A. Kissinger.
  \emph{Picturing Quantum Processes.}
  Cambridge University Press, 2017.

\bibitem{LeiferMilner} J. Leifer and R. Milner.
  \emph{Deriving bisimulation congruences for reactive systems.}
  In CONCUR, 2000.

\bibitem{SelingerValiron} P. Selinger and B. Valiron.
  \emph{A lambda calculus for quantum computation with classical control.}
  Mathematical Structures in Computer Science, 16(3), 2006.

\bibitem{Baez} J. Baez.
  \emph{The octonions.}
  Bulletin of the American Mathematical Society, 39(2), 2002.

\end{thebibliography}

\end{document}
